% ProveIt research continuation, 9 October 2026 (Pacific time).
% Replaces only the final rank-conjecture subsection of 05-signed-kernels.tex
% at commit 9bc738d3be22b8586a24693f19fb2e2a50ecd1bf.
% Full article, exact code and certificates accompany this fragment.

\subsection{Exact ranks from two polynomial reflections}
\label{l4rank:sec:rank}
Retain exactly the coefficient matrix $A_w$ defined above, including the
stuffle-minus-shuffle boundary rows and excluding the one divergent imaginary
coordinate. No distribution or higher-depth relations are added. Let
$A_w^{\neg g}$ denote the submatrix with the Gaussian columns deleted.
The following theorem replaces the former rank conjecture.

\begin{theorem}[All-weight ranks]
\label{signed:conj:rank}\label{l4rank:thm:rank}
For every odd $w\ge3$,
\begin{equation}\label{signed:eq:rank-conjecture}
 \operatorname{rank} A_w=5w-6,\qquad
 \operatorname{rank} A_w^{\neg g}=\frac{9w-11}{2}.
\end{equation}
For every even $w\ge2$,
\begin{equation}\label{l4rank:eq:even-rank}
 \operatorname{rank} A_w=6w-7,\qquad
 \operatorname{rank} A_w^{\neg g}=5w-6.
\end{equation}
In particular, the Gaussian-supported row space has dimension $(w-1)/2$
in odd weight and $w-1$ in even weight. In odd weight it is exactly the span
of the same-point shuffle rows.
\end{theorem}

\begin{proof}
Put $N=w-2$. Use six degree-$N$ polynomials $A,B,C,D,E,F$ to encode
a homogeneous null vector. Their respective colors are $(0,1),(1,0),(1,1),(1,2),(1,3),(2,1)$. Coefficients of $X^{a-1}Y^{b-1}$ have $a+b=w$.
Complete the missing coefficient of $A$ by
$v_{1,w-1}^{0,1}=-v_{w-1,1}^{1,0}$. This satisfies the missing stuffle row;
the missing shuffle row follows from its retained difference with stuffle.
The omitted coordinate occurs in no other row: outer index one in a
shuffle sum forces one of its factors to be $\Li_1(1)$.
Thus completion and restriction identify the matrix kernel with the full
homogeneous polynomial solution space.

Writing $P^{\mathsf T}(X,Y)=P(Y,X)$, the stuffle equations are
\begin{equation}\label{l4rank:eq:st}
 A=-B^{\mathsf T},\quad C=-C^{\mathsf T},\quad
 D=-F^{\mathsf T},\quad E=E^{\mathsf T}.
\end{equation}
The shuffle equations are
\begin{align}
 E(X+Y,X)+A(X+Y,Y)&=0,\notag\\
 B(X+Y,X)+B(X+Y,Y)&=0,\notag\\
 C(X+Y,Y)-F(X+Y,X)&=0,\notag\\
 D(X+Y,Y)-D(X+Y,X)&=0.\label{l4rank:eq:sh}
\end{align}
For general colors $(r,s)$ the shuffle polynomial is
$P_{s,r-s}(X+Y,X)+P_{r,s-r}(X+Y,Y)$; its coefficient of
$X^{p-1}Y^{q-1}$ is the stated binomial row. The four equations above come
from product colors $(0,1),(1,1),(1,2),(1,3)$. Swapping the factors,
conjugating, and discarding even-even imaginary rows exhaust the other cases.

The first two equations give
\[
 E(X,Y)=B(X-Y,X),\qquad B(X,Y)+B(X,X-Y)=0.
\]
With $\varepsilon=(-1)^N$, this implies
$E(Y,X)=-\varepsilon E(X,Y)$. Since $E=E^{\mathsf T}$, the whole $A,B,E$
block vanishes if $N$ is even. If $N=2m-1$ is odd, its free polynomial $B$
is precisely an anti-invariant under $Y\mapsto X-Y$, with basis
\begin{equation}\label{l4rank:eq:B-basis}
 B_j=X^{N-2j-1}(2Y-X)^{2j+1},\qquad 0\le j<m.
\end{equation}
Indeed, $U=X,V=2Y-X$ changes this reflection into $V\mapsto -V$.

The remaining equations give
\[
 F(X,Y)=C(X,X-Y),\quad D(X,Y)=-C(Y,Y-X),\quad C=-C^{\mathsf T}.
\]
They imply $D(X,X-Y)=-\varepsilon D(X,Y)$, while the last shuffle asks
for equality. Hence the $C,D,F$ block is zero when $N$ is even. When
$N=2m-1$, it has the basis determined by
\begin{equation}\label{l4rank:eq:C-basis}
 C_j=X^jY^{N-j}-X^{N-j}Y^j,\qquad 0\le j<m.
\end{equation}
Conversely, the displayed substitutions satisfy every homogeneous equation.
The $B_j$ and $C_j$ are independent and are retained in full after deleting
the missing coefficient of $A$. They therefore give an integral
$\mathbb Q$-basis of the full nullspace.

For odd $w=2m+1$, the nullity is $2m=w-1$, proving
$\operatorname{rank}A_w=(6w-7)-(w-1)=5w-6$. A null vector of
$A_w^{\neg g}$ extends by zero on the Gaussian columns; this is exactly
the $C$ block, of dimension $m$. Its $5w-6$ columns thus have rank
$5w-6-m=(9w-11)/2$. For even $w$ all columns of $A_w$ are independent,
and so are all columns of its submatrix.

Projection of the row space onto the non-Gaussian columns gives the stated
supported-row dimensions by taking the difference of the two ranks. In odd
weight $w=2m+1$, the same-point shuffle coefficient on $g_{a,w-a}$ is
$\binom{a-1}{p-1}+\binom{a-1}{w-p-1}$, $1\le p\le m$.
Its first $m$ columns form a triangular Pascal matrix with diagonal one.
These $m$ independent Gaussian-supported rows therefore exhaust that space.
\end{proof}

\paragraph{Exact membership and the harmonic-sum obstruction.}
Let $K_w$ have the preceding integer null vectors as columns. A rational
coefficient row $t$ is in the specified row space exactly when $tK_w=0$.
Otherwise any column with nonzero pairing is a separating certificate.
The basis is constructible in $O(w^2)$ integer arithmetic operations and
has $O(w)$-bit coefficients; coefficient recurrences and bit bounds are
given in the accompanying article.

For odd $w=2m+1$, choose $B=0$, $C=Y^{2m-1}-X^{2m-1}$.
It vanishes on every Gaussian column and has $D(1,0)=1$.
Since $S_{2m}=g_{2m,1}+\operatorname{Im}D_{2m,1}(\mathrm i,-1)$,
it pairs to one with $S_{2m}$. Thus no rational linear combination of these
fixed-weight product rows reduces $S_{2m}$ to Gaussian doubles and the
associated single-value right sides. In the formal quotient, adjoining it
to the Gaussian span increases the dimension from $m$ to $m+1$.
For the weight-five target above, this witness gives $T\cdot v=7$.
This is a uniform proof-system obstruction, not a numerical independence
theorem and compatible with the already proved $S_4$ identity, whose certificate uses a larger vocabulary.

\subsubsection{Restoring the right sides: an explicit compiler}
\label{l4rank:sec:compiler}
Set $\ell_k(r)=\Li_k(\mathrm i^r)$, with the formal regularization
$\ell_1(0)=0$, and define known degree-$N$ polynomials
\begin{align*}
 P_{r,s}(X,Y)&=\sum_{p+q=w}\operatorname{Im}(\ell_p(r)\ell_q(s))
                    X^{p-1}Y^{q-1},\\
 H_N(X,Y)&=\sum_{j=0}^{N}X^jY^{N-j},\\
 S_{01}&=P_{0,1}-\beta(w)H_N,\qquad
 S_{12}=P_{1,2}+\beta(w)H_N.
\end{align*}
The same six letters now encode the actual imaginary double values,
with the artificial missing coefficient $[X^0Y^N]A=-g_{w-1,1}-\beta(w)$.
Put
\begin{align*}
 e(X,Y)&=P_{0,1}(Y,X-Y)-S_{01}(X,X-Y),\\
 d(X,Y)&=S_{12}(X,Y)+P_{1,2}(X,Y-X),\\
 b(X,Y)&=P_{1,1}(Y,X-Y),\\
 h(X,Y)&=P_{1,3}(Y,Y-X)-e(Y,Y-X)+e(Y-X,Y).
\end{align*}

\begin{theorem}[Even-weight polynomial solution]\label{l4rank:thm:compiler}
For even $w\ge2$ the full solution is
\begin{align*}
 B(X,Y)&=\tfrac12\bigl(b(X,Y)+h(X,Y)\bigr),\\
 C(X,Y)&=\tfrac12\bigl(P_{1,3}(X,-Y)+P_{1,1}(Y,X)
                    -d(X-Y,-Y)+d(X-Y,X)\bigr),\\
 A(X,Y)&=S_{01}(X,Y)-B(Y,X),\\
 E(X,Y)&=B(X-Y,X)+e(X,Y),\\
 F(X,Y)&=C(X,X-Y)-P_{1,2}(Y,X-Y),\\
 D(X,Y)&=-C(Y,Y-X)+d(X,Y).
\end{align*}
Discarding the artificial coefficient gives explicit single-value reductions
for every convergent imaginary double at fourth roots of unity.
\end{theorem}
\begin{proof}
The affine stuffle equations are
$A+B^{\mathsf T}=S_{01}$, $C+C^{\mathsf T}=P_{1,1}$,
$D+F^{\mathsf T}=S_{12}$, $E-E^{\mathsf T}=P_{1,3}$.
The four shuffle right sides in \eqref{l4rank:eq:sh} are now
$P_{0,1},P_{1,1},P_{1,2},P_{1,3}$, respectively.
Elimination gives the last four displayed formulas and
$B(X,Y)+B(X,X-Y)=b(X,Y)$.
Even homogeneity and the $E$ stuffle give the difference
$B(X,Y)-B(X,X-Y)=h(X,Y)$, proving the formula for $B$.
For $C$, the last shuffle is
$D(u,u-v)-D(u,v)=P_{1,3}(v,u-v)$.
Use
$C(u-v,-v)=C(v-u,v)=P_{1,1}(v-u,v)-C(v,v-u)$
and then $(X,Y)=(v,v-u)$ to get the asserted formula.
Actual values satisfy these equations, and the zero homogeneous kernel
makes the solution unique.
\end{proof}

\paragraph{Two mixed leading-index families.}
Let $c_p=\operatorname{Im}D_{p,1}(\mathrm i,\mathrm i)$ and
$e_p=\operatorname{Im}D_{p,1}(\mathrm i,-\mathrm i)$. For all $m\ge1$,
\begin{align}
 c_{2m-1}&=(m-1)\beta(2m)-\tfrac12\log2\,\beta(2m-1)
 -\sum_{k=1}^{m-1}(1-2^{1-2k})\zeta(2k)\beta(2m-2k),
 \label{l4rank:eq:c-family}\\
 e_{2m-1}&=m\beta(2m)-\tfrac12\log2\,\beta(2m-1)
 -\sum_{k=1}^{m-1}\zeta(2k)\beta(2m-2k).
 \label{l4rank:eq:e-family}
\end{align}
For a direct coefficient proof, evaluate $C(1,0)$ and $E(1,0)$ in the compiler.
The resulting equations are
\begin{align*}
 2C(1,0)&=P_{1,3}(1,0)+P_{1,1}(0,1)-d(1,0)+d(1,1),\\
 2E(1,0)&=P_{1,1}(1,0)+P_{1,3}(1,0)+e(1,0)+e(0,1).
\end{align*}
In both lines the first two terms are $-\log2\,\beta(w-1)$, and
\begin{align*}
 d(1,1)-d(1,0)&=P_{1,2}(1,1)-P_{1,2}(1,-1)+N\beta(w),\\
 e(1,0)+e(0,1)&=-P_{0,1}(1,1)+P_{0,1}(1,-1)+w\beta(w).
\end{align*}
Only even inner, respectively even outer, indices survive the differences.
Using $\ell_{2k}(2)=-(1-2^{1-2k})\zeta(2k)$ and
$\operatorname{Im}\ell_j(1)=\beta(j)$ proves both formulas.

\paragraph{Status and evidence.}
The rank theorem proves both individual ranks, rather than inferring them
from their difference. The previous word ``Equivalently'' should therefore
be read as ``In particular.'' Independent matrix construction, exact integer
kernel checks and modular lower-bound certificates cover every weight
$2$--$41$; 210 even-weight formulas through weight 12 satisfy all 816 product
rows, and both leading-index formulas are checked through weight 16.
These finite checks support the uniform proofs; they are not their substitute.
The known general parity theorem already gives the relevant reduction
phenomenon. The new contribution here is the explicit solution of this
specified coefficient module. This row-system theorem asserts no numerical period independence. The
$S_4$ identity is proved separately in Theorem~\ref{gaussian:thm:S4}. Historical conjectural receipts are
retained as provenance rather than silently rewritten.
